\documentclass[aspectratio=43,10pt]{beamer}

%% ─── Theme ────────────────────────────────────────────────────────
\usetheme[iut, sans, progressbar=segmented, nosub]{Shibui}

%% ─── Packages ────────────────────────────────────────────────────
\usepackage[utf8]{inputenc}
\usepackage[T2A,T1]{fontenc}
\usepackage[french]{babel}
\usepackage{graphicx}
\usepackage{mathtools}
\usepackage{tikz}
\usepackage{pgfplots}
\pgfplotsset{compat=1.18}
\usetikzlibrary{arrows.meta,positioning,calc}

\usepackage{amsmath,amsfonts,bm}

%% ─── Operators ───────────────────────────────────────────────────
\DeclareMathOperator*{\argmax}{arg\,max}
\DeclareMathOperator*{\argmin}{arg\,min}
\DeclareMathOperator{\MSE}{MSE}
\DeclareMathOperator{\sign}{sign}
\DeclareMathOperator{\sinc}{sinc}
\DeclareMathOperator{\tr}{tr}

%% ─── Number sets ─────────────────────────────────────────────────
\def\R{{\mathbb{R}}}
\def\C{{\mathbb{C}}}
\def\N{{\mathbb{N}}}
\def\Z{{\mathbb{Z}}}

%% ─── Infinitesimal and units ────────────────────────────────────
\def\d{{\rm d}}
\def\Hz{{\mbox{Hz}}}
\def\m{{\mbox{m}}}

%% ─── Complex / signal ───────────────────────────────────────────
\DeclareMathOperator{\complexj}{\mathsf{j}}   % imaginary unit

%% ─── Probability ─────────────────────────────────────────────────
\newcommand{\Esp}[2]{\mathbb{E}_{#1}\!\left[#2\right]}
\renewcommand{\Pr}{\mathbb{P}}

%% ─── Linear algebra ──────────────────────────────────────────────
\newcommand{\abs}[1]{\left|{#1}\right|}
\newcommand{\norm}[1]{\left\|#1\right\|}
\newcommand{\inner}[2]{\left\langle#1,\,#2\right\rangle}

%% ─── Logic ───────────────────────────────────────────────────────
\newcommand{\then}{\Rightarrow}
\newcommand{\iif}{\Leftrightarrow}

%% ─── Draft annotation ────────────────────────────────────────────
\newcommand{\red}[1]{{\color{red}#1}}

%% ─── Shared figure color (used by common/figures/*/*.tex bodies) ──
\definecolor{niceblue}{RGB}{0,128,172}

%% ─── Rappel header (styled subtitle for reference-card asides) ───
%% Requires \definecolor{niceblue}{...} to be defined by the including document.
\newcommand{\rappelhead}[1]{%
    \bigskip
    {\color{niceblue}\bfseries #1}\\[-0.4em]
    {\color{niceblue!35}\rule{\linewidth}{0.6pt}}
    \smallskip

}


%% ─── Metadata ─────────────────────────────────────────────────────
\title{OML3}
\subtitle{Chapitre 3 \\ La Transformée de Fourier Discrète}
\author{\textbf{BUT2} -- Génie Électrique et Informatique Industrielle}
\date{2025-2026}

%% ─── Theorem environments ─────────────────────────────────────────
\newtheorem{proposition}{Propriété}
\newtheorem{remark}{Remarque}

%% ─── Text helpers (mirrors poly/main.tex) ──────────────────────────
\newcommand{\highlight}[1]{{\color{niceblue}\emph{#1}}}
\newcommand{\tf}{\overset{\mathcal F}{\longleftrightarrow}}
\definecolor{niceblue}{RGB}{0,128,172}

%% ─── Peigne de Dirac (Ш) ────────────────────────────────────────────
%% Real Cyrillic Sha via T2A encoding (cm-super provides the outline),
%% same approach as elec-s3/ch2 -- confirmed working here too.
\newcommand{\Sha}{\text{\fontencoding{T2A}\selectfont\CYRSH}}

%% ─── Placeholder for live-tool captures (Ableton/Audacity) ─────────
\newcommand{\captureplaceholder}[1]{%
    \begin{center}
    \fbox{\parbox{0.75\linewidth}{\centering\vspace{1.5em}\textit{[CAPTURE À VENIR : #1]}\vspace{1.5em}}}
    \end{center}
}

%% ─── Beat graph styles (rappel-graph*.tex figures rely on these) ───
%% Native Shibui (iut palette) colors, not niceblue : shibuidarkgray is
%% the theme's blue, used as the neutral box/arrow color ; red is kept
%% for its own sake, reserved for actually highlighting the "current
%% position" node (beatfilled), not used anywhere else.
\tikzset{
    beatbox/.style={draw=shibuidarkgray, thick, rounded corners, minimum height=1.4cm,
                text width=2.4cm, align=center, font=\small},
    beatfilled/.style={beatbox, draw=red, fill=white!85!red},
    beatarr/.style={-{Latex[length=2.5mm]}, shibuidarkgray, thick}
}

%% ─── Figure helper : fit any \input'd figure to a target width ────
\newcommand{\slidefig}[2]{%
\begin{center}
\resizebox{#1}{!}{\input{../../../common/figures/ch3/#2}}
\end{center}
}

%% ─────────────────────────────────────────────────────────────────
\begin{document}

\maketitle

%% ═══════════════════════════════════════════════════════════════
\section{Rappel}

\begin{frame}{Décomposition spectrale}
    \slidefig{0.65\textwidth}{decompose-porte-fourier.tex}
\end{frame}

\begin{frame}{}
    \slidefig{\textwidth}{rappel-graph.tex}
    \vspace{1em}
    % \begin{center}
    %     \small Aujourd'hui : discret $\Leftrightarrow$ périodique, appliqué deux fois.
    % \end{center}
\end{frame}

\begin{frame}{Séries de Fourier}
    \begin{center}
        \small fonctions \highlight{périodiques} (période $T$)
    \end{center}
    \vspace{0.8em}
    \begin{center}
        \resizebox{\textwidth}{!}{%
        $x(t) = \displaystyle\sum_{n=-\infty}^{+\infty} c_n\,e^{2\complexj\pi nt/T}, \qquad
          c_n = \dfrac1T\displaystyle\int_T x(t)\,e^{-2\complexj\pi nt/T}\,\mathrm dt$%
        }
    \end{center}
\end{frame}

\begin{frame}{Transformée de Fourier}
    \begin{center}
        \small fonctions \highlight{non périodiques}
    \end{center}
    \vspace{0.8em}
    \begin{center}
        \Large
        $$X(f) = \int_{-\infty}^{+\infty} x(t)\,e^{-2\complexj\pi ft}\,\mathrm dt$$
    \end{center}
\end{frame}

\begin{frame}{L'impulsion de Dirac}
    \slidefig{0.55\textwidth}{dirac-sifting.tex}
    \begin{center}
        $$\int_{-\infty}^{+\infty} x(t)\,\delta(t-t_0)\,\mathrm dt = x(t_0)$$
    \end{center}
\end{frame}

\begin{frame}{Produit de convolution}
    \begin{center}
        \Large
        $$(x*y)(t) = \int_{-\infty}^{+\infty} x(\tau)\,y(t-\tau)\,\mathrm d\tau$$
    \end{center}
\end{frame}

\begin{frame}{Dualité produit / convolution}
    \begin{center}
        \Large
        $$x(t)\,y(t) \;\tf\; (X*Y)(f)$$
        \vspace{0.5em}
        $$(x*y)(t) \;\tf\; X(f)\,Y(f)$$
    \end{center}
\end{frame}

%% ═══════════════════════════════════════════════════════════════
\section{Fonction discrète}


\begin{frame}{Définition}
    \begin{center}
        \Large
        $$x:\Z\to\R,\qquad n\mapsto x[n]$$
    \end{center}
\end{frame}

\begin{frame}{Deux exemples de fonctions discrètes}
    \begin{columns}
        \column{0.5\textwidth}
        \centering
        \textbf{Échelon unité}
        $$u[n] = \begin{cases} 1 & n\geq0 \\ 0 & n<0 \end{cases}$$
        \column{0.5\textwidth}
        \centering
        \textbf{Impulsion unité}
        $$\delta[n] = \begin{cases} 1 & n=0 \\ 0 & n\neq0 \end{cases}$$
    \end{columns}
\end{frame}

%% ═══════════════════════════════════════════════════════════════
\section{Discret \texorpdfstring{$\Leftrightarrow$}{<=>} périodique}

\begin{frame}{Le peigne de Dirac}
    \slidefig{0.6\textwidth}{peigne-dirac.tex}
    \begin{center}
        \large
        $$\Sha_{T_e}(t) = \sum_{n=-\infty}^{+\infty}\delta(t-nT_e)$$
    \end{center}
\end{frame}

\begin{frame}{À connaître : La transformée de Fourier d'un peigne de Dirac}

     \begin{center}
        \Large
        $$\Sha_{T_e}(t) = \sum_{n=-\infty}^{+\infty}\delta(t-nT_e)$$
    \end{center}
    \vspace{2em}
    \begin{center}
        \Large
        $$\Sha_{T_e}(t) \;\tf\; F_e\,\Sha_{F_e}(f), \qquad F_e=1/T_e$$
    \end{center}
    \vspace{1em}
\end{frame}

\begin{frame}{$x$ discret $\Rightarrow X_e$ périodique}
    \begin{columns}
        \column{0.5\textwidth}
        \centering
        \slidefig{\textwidth}{echantillonnage-temps.tex}
        \begin{center}\small temps : $x_e$ discret\end{center}
        \onslide<2->{%
        $$
        \begin{aligned}
            x_e(t) &= x(t)\times\Sha_{T_e}(t)  \\ &= \sum_{n=-\infty}^{+\infty} x(nT_e)\, \delta(t-nT_e)
        \end{aligned}
        $$
        }

        \column{0.5\textwidth}
        \centering
        \slidefig{\textwidth}{echantillonnage-frequence.tex}
        \begin{center}\small fréquence : $X_e$ périodique\end{center}
        \onslide<2->{%
        $$
        \begin{aligned}
            X_e(f) &= \int_{-\infty}^{+\infty} x_e(t)\, e^{-\complexj 2\pi ft}\,\mathrm dt  \\ &= \sum_{n=-\infty}^{+\infty} x(nT_e)\, e^{-\complexj 2\pi f nT_e}
        \end{aligned}
        $$
        }
    \end{columns}
\end{frame}

\begin{frame}{$x$ périodique $\Rightarrow X$ discret}
    \begin{columns}
        \column{0.5\textwidth}
        \centering
        \slidefig{\textwidth}{periodique-temps.tex}
        \begin{center}\small temps : $x$ périodique\end{center}
        \onslide<2->{%
        $$
        \begin{aligned}
            x(t) &= x(t+T) \\ &= \sum_{n=-\infty}^{+\infty} c_n\, e^{2\complexj\pi nt/T}
        \end{aligned}
        $$
        }

        \column{0.5\textwidth}
        \centering
        \slidefig{\textwidth}{periodique-frequence.tex}
        \begin{center}\small fréquence : $X$ discret\end{center}
        \onslide<2->{%
        $$
        \begin{aligned}
            X(f) &= \mathcal F\Big[\sum_{n=-\infty}^{+\infty} c_n\, e^{2\complexj\pi nt/T}\Big](f) \\ &= \sum_{n=-\infty}^{+\infty} c_n\, \delta(f-n/T)
        \end{aligned}
        $$
        }
    \end{columns}
\end{frame}

%% ═══════════════════════════════════════════════════════════════
% \section{Bilan (séance 1)}

\begin{frame}{Où on en est}
    \slidefig{\textwidth}{rappel-graph-mid.tex}
    \vspace{1em}
    \begin{center}
        Discret $\Leftrightarrow$ périodique : établi dans les deux sens.
    \end{center}
\end{frame}

\begin{frame}{Fréquence normalisée}


    % \begin{center}
        $$
            \Large
            \begin{aligned}
            X_e({\color{shibuidarkgray}f}) &= \sum_{n=-\infty}^{+\infty} \red{x(nT_e)}\, e^{-\complexj 2\pi n {\color{red}fT_e}} 
            \\
             X(\red{\nu}) &= \sum_{n=-\infty}^{+\infty} \red{x[n]}\,e^{-2\complexj\pi\red{\nu} n}
            \end{aligned}
        $$

    % \end{center}

    \begin{center}
        \Large
        $$(t,f) \;\longleftrightarrow\; (n,\nu)$$
    \end{center}
    \begin{center}
        \small temps continu / indice discret \hspace{3em} fréquence / fréquence normalisée
    \end{center}
    \vspace{1.5em}

\end{frame}

\begin{frame}{Presque...}
    \begin{columns}
        \column{0.5\textwidth}
        \centering
        \slidefig{\textwidth}{dtft-synthetique-temps.tex}
        \begin{center}\small temps : $x[n]$ fini (non périodique)\end{center}
        \column{0.5\textwidth}
        \centering
        \slidefig{\textwidth}{dtft-synthetique-frequence.tex}
        \begin{center}\small fréquence : $X(\nu)$ périodique... mais continue\end{center}
    \end{columns}
    \vspace{0.5em}
    \begin{center}
        Support borné $\nu\in[0,1)$ -- mais pas discret.\\[0.3em]
        \Large Exploitable numériquement ?
    \end{center}
\end{frame}

%% ═══════════════════════════════════════════════════════════════
\section{La TFD}

% \begin{frame}{Fréquence normalisée}
%     \begin{center}
%         \Large
%         $$\nu = f/F_e$$
%     \end{center}
%     \vspace{0.5em}
%     \begin{center}
%         $$X(\nu) = \sum_{n=-\infty}^{+\infty} x[n]\,e^{-2\complexj\pi\nu n}, \qquad \text{période } 1$$
%     \end{center}
% \end{frame}

\begin{frame}{La Transformée de Fourier Discrète (TFD)}
    \begin{center}
        \Large
        $$X[k] = \sum_{n=0}^{N-1} x[n]\,e^{-2\complexj\pi kn/N}, \qquad k=0,\dots,N-1$$
    \end{center}
\end{frame}

\begin{frame}{La TFD inverse}
    \begin{center}
        \Large
        $$x[n] = \frac1N\sum_{k=0}^{N-1} X[k]\,e^{2\complexj\pi kn/N}, \qquad n=0,\dots,N-1$$
    \end{center}
\end{frame}

% \begin{frame}{Pourquoi c'est important}
%     \captureplaceholder{écran d'un analyseur de spectre / oscilloscope en mode FFT}
%     \begin{center}
%         \small C'est ce calcul, partout.
%     \end{center}
% \end{frame}



\end{document}
